% Source-faithful mathematical blueprint of Hartshorne, Algebraic Geometry (1977), Appendix A.
% Authorial exposition, examples, and exercises are omitted; mathematical conventions, statements,
% formulas, and proof arguments are retained.
\section{Intersection Theory}
\label{ha-appA-sec1}
\source{hartshorne-algebraic-geometry:A.1}

\begin{definition}[Cycles, rational equivalence, and Chow groups]
\source{hartshorne-algebraic-geometry:A.1}
For a variety $X$ over an algebraically closed field, a codimension-$r$ cycle
is a finite integral combination of closed irreducible codimension-$r$
subvarieties.  A closed subscheme of pure codimension $r$ determines the
cycle whose coefficient at a component is the length at its generic point.
For $f:X\to X'$, proper pushforward sends a subvariety $Y$ to zero if its
dimension drops and otherwise to $[K(Y):K(f(Y))]f(Y)$, extended linearly.
If $\widetilde V\to V$ is the normalization, linearly equivalent Weil
divisors on $\widetilde V$ define rationally equivalent cycles on $X$.
Write $A^r(X)$ for codimension-$r$ cycles modulo this relation and
$A(X)=\bigoplus_rA^r(X)$.
\end{definition}

An intersection theory on nonsingular quasi-projective varieties has products
$A^r(X)\times A^s(X)\to A^{r+s}(X)$ and pullbacks defined by the graph, subject
to the following axioms.

\begin{convention}[Chow-ring structure]
\label{ha-appA-A1}
\source{hartshorne-algebraic-geometry:A1}
\dcref{A1}
$A(X)$ is a commutative associative graded ring with identity.
\end{convention}

\begin{convention}[Functorial pullback]
\label{ha-appA-A2}
\source{hartshorne-algebraic-geometry:A2}
\dcref{A2}
Pullback is a ring map and $(gf)^*=f^*g^*$.
\end{convention}

\begin{convention}[Functorial proper pushforward]
\label{ha-appA-A3}
\source{hartshorne-algebraic-geometry:A3}
\dcref{A3}
Proper pushforward is graded and $(gf)_*=g_*f_*$.
\end{convention}

\begin{convention}[Projection formula]
\label{ha-appA-A4}
\source{hartshorne-algebraic-geometry:A4}
\dcref{A4}
$f_*(x f^*y)=f_*(x)y$.
\end{convention}

\begin{convention}[Diagonal formula]
\label{ha-appA-A5}
\source{hartshorne-algebraic-geometry:A5}
\dcref{A5}
For the diagonal $\Delta:X\to X\times X$, $YZ=\Delta^*(Y\times Z)$.
\end{convention}

\begin{convention}[Proper-intersection multiplicities]
\label{ha-appA-A6}
\source{hartshorne-algebraic-geometry:A6}
\dcref{A6}
If $Y,Z$ meet properly, $YZ=\sum_W i(Y,Z;W)W$, with $i(Y,Z;W)$ local at the
generic point of $W$.
\end{convention}

\begin{convention}[Cartier-divisor normalization]
\label{ha-appA-A7}
\source{hartshorne-algebraic-geometry:A7}
\dcref{A7}
If $Z$ is an effective Cartier divisor meeting $Y$ properly, $YZ$ is the cycle
of the Cartier divisor cut on $Y$.
\end{convention}

\begin{theorem}[Intersection theory]
\label{ha-appA-thm1-1}
\dcref{A.1.1}
On nonsingular quasi-projective varieties there is a unique intersection
theory satisfying A1--A7.  For proper intersection, the local multiplicity is
\[
i(Y,Z;W)=\sum_{j\ge0}(-1)^j\length_A\operatorname{Tor}^A_j(A/\mathfrak a,A/\mathfrak b),
\]
where $A=\mathcal O_{X,\eta_W}$ and $\mathfrak a,\mathfrak b$ are the ideals
of $Y,Z$.  Chow's moving lemma says that $Z$ is rationally equivalent to a
cycle meeting $Y$ properly, and two such choices give the same product.
\source{hartshorne-algebraic-geometry:A.1.1}
\uses{ha-appA-A1,ha-appA-A2,ha-appA-A3,ha-appA-A4,ha-appA-A5,ha-appA-A6,ha-appA-A7}
\end{theorem}

\begin{proof}
Serre's Tor formula has the required local, normalization, and positivity
properties.  By the moving lemma replace one cycle by a rationally equivalent
cycle meeting the other properly.  For uniqueness use A5 to reduce to the
diagonal in $X\times X'$, which is a local complete intersection; A6 reduces
to one factor a complete intersection of Cartier divisors, and repeated A7
then determines the product.  This also proves independence of the moved
representative.
\end{proof}



\section{Properties of the Chow Ring}
\label{ha-appA-sec2}
\source{hartshorne-algebraic-geometry:A.2}

\begin{convention}[Divisors and the Picard group]
\label{ha-appA-A8}
\source{hartshorne-algebraic-geometry:A8}
\dcref{A8}
$A^1(X)\simeq\operatorname{Pic}X$.
\end{convention}

\begin{convention}[Affine-space invariance]
\label{ha-appA-A9}
\source{hartshorne-algebraic-geometry:A9}
\dcref{A9}
Projection $X\times\mathbb A^n\to X$ induces
$A(X)\simeq A(X\times\mathbb A^n)$.
\end{convention}

\begin{convention}[Localization]
\label{ha-appA-A10}
\source{hartshorne-algebraic-geometry:A10}
\dcref{A10}
For a nonsingular closed $Y\subset X$, with $U=X-Y$, the sequence
$A(Y)\to A(X)\to A(U)\to0$ is exact.
\end{convention}

\begin{convention}[Projective-bundle formula]
\label{ha-appA-A11}
\source{hartshorne-algebraic-geometry:A11}
\dcref{A11}
If $\mathcal E$ has rank $r$, then $A(\mathbb P(\mathcal E))$ is a free
$A(X)$-module on $1,\xi,\ldots,\xi^{r-1}$, where $\xi$ is the tautological
divisor class.
\end{convention}

\section{Chern Classes}
\label{ha-appA-sec3}
\source{hartshorne-algebraic-geometry:A.3}

\begin{definition}[Chern classes]
\source{hartshorne-algebraic-geometry:A.3}
For a rank-$r$ locally free sheaf $\mathcal E$ on $X$, define
$c_i(\mathcal E)\in A^i(X)$ by $c_0=1$ and
\[
\sum_{i=0}^r(-1)^i\xi^{r-i}\pi^*c_i(\mathcal E)=0
\quad\text{in }A(\mathbb P(\mathcal E)),
\]
where $\pi$ is the projection.  The total class is
$c(\mathcal E)=\sum_i c_i(\mathcal E)$.
\end{definition}

\begin{convention}[First Chern class of a line bundle]
\label{ha-appA-C1}
\source{hartshorne-algebraic-geometry:C1}
\dcref{C1}
$c(\mathcal O_X(D))=1+D$.
\end{convention}

\begin{convention}[Pullback of Chern classes]
\label{ha-appA-C2}
\source{hartshorne-algebraic-geometry:C2}
\dcref{C2}
$c_i(f^*\mathcal E)=f^*c_i(\mathcal E)$.
\end{convention}

\begin{convention}[Whitney sum formula]
\label{ha-appA-C3}
\source{hartshorne-algebraic-geometry:C3}
\dcref{C3}
For $0\to\mathcal E'\to\mathcal E\to\mathcal E''\to0$,
$c(\mathcal E)=c(\mathcal E')c(\mathcal E'')$.
\end{convention}

\begin{convention}[Splitting principle]
\label{ha-appA-C4}
\source{hartshorne-algebraic-geometry:C4}
\dcref{C4}
If $\mathcal E$ has a filtration with invertible quotients $\mathcal L_i$,
then $c(\mathcal E)=\prod_i(1+c_1(\mathcal L_i))$.
\end{convention}

\begin{convention}[Chern-root identities]
\label{ha-appA-C5}
\source{hartshorne-algebraic-geometry:C5}
\dcref{C5}
Under formal Chern roots $a_i,b_j$,
$c(\mathcal E\otimes\mathcal F)=\prod_{i,j}(1+(a_i+b_j)t)$,
$c(\bigwedge^p\mathcal E)=\prod_{|I|=p}(1+(\sum_{i\in I}a_i)t)$, and
$c(\mathcal E^\vee)=\prod_i(1-a_it)$.
\end{convention}

\begin{convention}[Zero scheme of a section]
\label{ha-appA-C6}
\source{hartshorne-algebraic-geometry:C6}
\dcref{C6}
If a section of $\mathcal E$ has zero scheme $Y$ of codimension $r$, then
$c_r(\mathcal E)=[Y]$.
\end{convention}

\begin{convention}[Self-intersection formula]
\label{ha-appA-C7}
\source{hartshorne-algebraic-geometry:C7}
\dcref{C7}
If $i:Y\hookrightarrow X$ is nonsingular of codimension $r$, then
$i_*c_r(N_{Y/X})=Y\cdot Y$.
\end{convention}

\section{The Riemann--Roch Theorem}
\label{ha-appA-sec4}
\source{hartshorne-algebraic-geometry:A.4}

\begin{definition}[Chern character and Todd class]
\source{hartshorne-algebraic-geometry:A.4}
For formal roots $a_i$ of $\mathcal E$, set
$\operatorname{ch}(\mathcal E)=\sum_i e^{a_i}$ and
$\operatorname{td}(\mathcal E)=\prod_i a_i/(1-e^{-a_i})$.  Thus
\[
\operatorname{ch}=r+c_1+\tfrac12(c_1^2-2c_2)+\tfrac16(c_1^3-3c_1c_2+3c_3)+\cdots.
\]
\end{definition}

\begin{theorem}[Hirzebruch--Riemann--Roch]
\label{ha-appA-thm4-1}
\dcref{A.4.1}
If $\mathcal E$ is locally free on a nonsingular projective $n$-fold $X$,
\[
\chi(\mathcal E)=\deg\bigl(\operatorname{ch}(\mathcal E)\operatorname{td}(T_X)\bigr)_n.
\]
\source{hartshorne-algebraic-geometry:A.4.1}
\end{theorem}
\begin{proof}
The source attributes this theorem to Hirzebruch over $\mathbb C$ and to
Grothendieck in generalized form; it does not give a proof here.
\end{proof}

\section{Complements and Generalizations}
\label{ha-appA-sec5}
\source{hartshorne-algebraic-geometry:A.5}

\begin{theorem}[Nakai--Moishezon criterion]
\label{ha-appA-thm5-1}
\dcref{A.5.1}
For a Cartier divisor $D$ on a scheme $X$ proper over an algebraically closed
field, $D$ is ample iff
$D^{\dim Y}\cdot Y>0$ for every positive-dimensional integral closed
subvariety $Y\subset X$.
\source{hartshorne-algebraic-geometry:A.5.1}
\end{theorem}
\begin{proof}
The source attributes this criterion to Nakai and says that its proof is
clarified and simplified by Kleiman; no proof is supplied in this appendix.
\end{proof}

\begin{theorem}[Hodge index theorem]
\label{ha-appA-thm5-2}
\dcref{A.5.2}
If $X$ is a nonsingular projective surface, $H$ is ample, and $D$ is a
numerically nonzero divisor, then
\[
(D\cdot H)^2\geq (H^2)(D^2),
\]
with equality if and only if $D$ is numerically proportional to $H$.
\source{hartshorne-algebraic-geometry:A.5.2}
\end{theorem}
\begin{proof}
The source attributes this theorem to Hodge and says that the proof uses
Hodge's theory of harmonic integrals; no proof is supplied here.
\end{proof}

\begin{definition}[K-theory pushforward]
\source{hartshorne-algebraic-geometry:A.5}
The Chern classes extend additively to $K(X)$, the Chern character is a ring
map $K(X)\to A(X)\otimes\mathbb Q$, and a proper morphism has
$f_![\mathcal F]=\sum_i(-1)^i[R^if_*\mathcal F]$.
\end{definition}

\begin{theorem}[Grothendieck--Riemann--Roch]
\label{ha-appA-thm5-3}
\dcref{A.5.3}
For a smooth projective morphism $f:X\to Y$ of nonsingular quasi-projective
varieties and $x\in K(X)$,
\[
\operatorname{ch}(f_!x)\operatorname{td}(T_Y)
=f_*\bigl(\operatorname{ch}(x)\operatorname{td}(T_X)\bigr).
\]
\source{hartshorne-algebraic-geometry:A.5.3}
\end{theorem}
\begin{proof}
The source attributes this theorem to Grothendieck and refers to its proof in
the cited literature; no proof is supplied here.
\end{proof}
